\subsection{One fixed finite type and host for a full weight progression}
\label{sec:cubic-uniform-weight-type}

This section retains one completed instance of AC-G78--AC-G83 at
\(p=269\), including the separated residual characters, the fixed
CM field \(F\), the actual auxiliary prime \(q\), and every old,
auxiliary-neat, and coefficient-prime projector.
The finite group at the selected \(q\)-place is
\(\Gamma_q=\operatorname{GL}_3(\mathbb F_q)\).
Let \(\tau=\tau_\theta\) be the irreducible characteristic-zero
finite representation constructed in AC-G72-character, where
\(\theta\) has order \(269\).
Enlarge the coefficient field once to define \(\tau\), all fixed
characters, and the embeddings used for weights, and choose a
\(\Gamma_q\)-stable finite free lattice \(\tau_{\mathcal O}\).
Such a lattice is obtained by taking the \(\mathcal O\)-span of the
finite group orbit of a basis.

Write
\[
 T=\mathfrak t_p,\qquad T_d=\mathfrak t_p(d),\qquad
 D_d=T/T_d,\qquad
 \Lambda_p=\mathcal O[[T]],\qquad A_d=\mathcal O[D_d\times\Gamma_q].
 \tag{G87.1}
\]
Here \(T=T_1\) is the pro-\(p\) torus in the literal source
ordinary construction.  In the retained field \(F\), its local
fields at \(p\) are unramified, so \(T\) is a finite product of
copies of \(\mathbb Z_p\) under the \(p\)-adic logarithm.
Thus \(\Lambda_p\) is a complete local power-series ring.
The larger coefficient field is used to state automorphic occurrence;
no original deformation coefficient hypothesis is silently changed.

\begin{lemma}[AC-G87-projective: the actual tower with its finite group action]
\label{lem:G87-projective}
At the fixed full principal \(q\)-level, let \(P_d\) be the weight-zero
finite-level ordinary module with all of the fixed projectors
\(e_{\rm aux},e_{\rm old},e_{\delta}\) imposed from the start and
localized at the fixed residual maximal ideal \(\mathfrak m\)
of the away Hecke algebra.
Then \(P_d\) is a finite projective \(A_d\)-module.
The ordinary trace transition induces
\[
 A_d\otimes_{A_{d+1}}P_{d+1}\ \simeq\ P_d.
 \tag{G87.2}
\]
Consequently the inverse-trace limit
\[
 P=\varprojlim_d P_d
 \tag{G87.3}
\]
is finite projective over \(\Lambda_p[\Gamma_q]\).

For every continuous weight
\(\lambda:T\to\mathcal O'{}^\times\), where \(\mathcal O'\) is
the ring of integers of a finite coefficient extension, its actual
specialization is
\[
 P\otimes_{\Lambda_p,\,w_0\lambda^{-1}}\mathcal O'
 \ \simeq\
 e_{\rm aux}e_{\rm old}e_\delta
 H_\lambda(K^{p,q}K_q^1)_{\mathfrak m}.
 \tag{G87.4}
\]
This isomorphism respects \(\Gamma_q\) and the away Hecke operators.
\end{lemma}

\begin{proof}
\noindent\textbf{The simultaneous finite deck group.}
Begin at the sufficiently small auxiliary level used in AC-G72,
before the retained auxiliary projector is applied.
At \(q\), use \(K_q^1=\ker(\operatorname{GL}_3(\mathcal O_q)
\to\Gamma_q)\).
At \(p\), the quotient between the levels
\(\operatorname{Iw}(d,d)\) and \(\operatorname{Iw}(1,d)\)
is \(D_d\).
These groups act at different places and commute.
The raw weight-zero module is the module of functions on a finite
double-coset set.  The action of \(D_d\times\Gamma_q\) on that set
is free: a stabilizer would be a rational element lying in a conjugate
of the larger compact level, and the sufficiently small away level
kills that finite stabilizer.
One can therefore choose one point in each orbit; the raw module is
a finite free \(\mathcal O[D_d\times\Gamma_q]\)-module.

The ordinary projector is an integral direct-summand projector by
Thorne's Lemma 6.1.  It commutes with both deck actions.
The three already constructed projectors
\(e_{\rm aux},e_{\rm old},e_\delta\) are integral, commute with the
same actions, and are imposed before any eigensystem is selected.
Localization at the fixed away-Hecke maximal ideal is also an
integral direct summand on these finite \(\mathcal O\)-modules:
the finite image Hecke algebra is a finite algebra over the complete
DVR \(\mathcal O\), and decomposes into its local factors.
Equivalently one can localize at the fixed multiplicative set in
the abstract away Hecke algebra.
All these operations commute with the maps in question.
Their common image \(P_d\) is thus projective over \(A_d\).
There is no averaging over \(|\Gamma_q|\).

\noindent\textbf{The transition.}
At the common lower-unipotent depth \(d+1\), Thorne's equation (6.1)
identifies the torus coinvariants with the lower torus level by trace.
The ordinary depth-stabilization immediately following (6.1) identifies
that lower torus level with depth \(d\).
All maps commute with the action at \(q\), the away Hecke algebra,
and the retained projectors.
This proves (G87.2).  It also proves surjectivity of the transition.
The relation between inclusion and trace is the literal norm identity
\[
 \operatorname{incl}\bigl(\operatorname{Tr}(v)\bigr)
   =\sum_{a\in\ker(D_{d+1}\to D_d)}a v.
 \tag{G87.5}
\]
No division by this \(p\)-power group order occurs.

\noindent\textbf{Projectivity of the inverse limit.}
Since \(D_1=1\), the module \(P_1\) is finite projective over
\(\mathcal O[\Gamma_q]\).
Choose an \(\mathcal O[\Gamma_q]\)-linear idempotent presenting it
as a summand of a finite free module.
Use the same constant idempotent to define a projective
\(\Lambda_p[\Gamma_q]\)-module \(Q\), and put
\(Q_d=A_d\otimes_{\Lambda_p[\Gamma_q]}Q\).
We construct compatible isomorphisms \(Q_d\simeq P_d\).
The first is the chosen presentation.
If an isomorphism has been constructed at depth \(d\), projectivity
of \(Q_{d+1}\) lifts the resulting map to \(P_{d+1}\).
The kernel of \(A_{d+1}\to A_d\) lies in the Jacobson radical:
modulo the uniformizer it is the augmentation ideal of a finite
\(p\)-group, and is nilpotent.  Nakayama therefore makes the lift
surjective. Since \(P_{d+1}\) is projective, the surjection splits;
its kernel has zero reduction to \(A_d\), so Nakayama makes the kernel
zero. This produces the required compatible isomorphisms.
Finally
\(\varprojlim_d A_d=\Lambda_p[\Gamma_q]\);
passing to the images of the constant idempotent gives
\(P\simeq Q\).  This proves finite projectivity without assuming that
an arbitrary inverse limit of projective modules is projective.

\noindent\textbf{Actual weight specialization.}
Modulo any fixed coefficient power, choose \(d\) so that the
continuous character \(w_0\lambda^{-1}\) factors through \(D_d\).
On a regular finite group-algebra module the twisted norm
\[
 [v]\ \longmapsto\
 \sum_{a\in D_d}(w_0\lambda^{-1})(a)^{-1}a v
 \tag{G87.6}
\]
identifies character coinvariants with the corresponding eigenspace.
It is an isomorphism by the coefficient-at-identity calculation,
not by division by \(|D_d|\).
The same statement holds on a projective summand.
It is natural for every operator commuting with \(D_d\), in
particular for \(\Gamma_q\) and all the retained projectors.
The norm maps compose with the trace transitions by (G87.5).
The finite-coefficient comparison in the complete proof of
Thorne's Proposition 6.2 identifies these eigenspaces with the
reductions of the actual \(H_\lambda\).
Taking the coefficient inverse limit proves (G87.4).
This is precisely the actual-module control argument of
AC-G78-module, now retaining the finite group at \(q\).
\end{proof}

\begin{proposition}[AC-G87-type: finite-type multiplicity commutes with every weight]
\label{prop:G87-type}
The module
\[
 \mathcal M_\tau
  =\operatorname{Hom}_{\mathcal O[\Gamma_q]}
       (\tau_{\mathcal O},P)
 \tag{G87.7}
\]
is finite free over \(\Lambda_p\).
For every commutative \(\Lambda_p\)-algebra \(R\), the natural map
\[
 R\otimes_{\Lambda_p}\mathcal M_\tau
 \ \longrightarrow\
 \operatorname{Hom}_{R[\Gamma_q]}
  (R\otimes_{\mathcal O}\tau_{\mathcal O},
       R\otimes_{\Lambda_p}P)
 \tag{G87.8}
\]
is an isomorphism.
If \(\tau\) occurs in one characteristic-zero weight fiber of \(P\),
then \(\mathcal M_\tau\) has positive rank, and \(\tau\) occurs
in every characteristic-zero weight fiber.
\end{proposition}

\begin{proof}
For any commutative \(\mathcal O\)-algebra \(R\), there is a natural
isomorphism
\[
 \operatorname{Hom}_{R[\Gamma_q]}
       (R\otimes_{\mathcal O}\tau_{\mathcal O},R[\Gamma_q])
 \ \simeq\
 \operatorname{Hom}_{\mathcal O}(\tau_{\mathcal O},R)
 \ \simeq\ R\otimes_{\mathcal O}\tau_{\mathcal O}^{\vee}.
 \tag{G87.9}
\]
The first map takes the coefficient of the identity group element.
Its inverse sends a linear functional \(\ell\) to
\[
 v\ \longmapsto\ \sum_{g\in\Gamma_q}\ell(g^{-1}v)\,g.
 \tag{G87.10}
\]
The displayed formula is equivariant, and direct substitution checks
that the two maps are inverse.
Since \(\tau_{\mathcal O}\) is finite free, the second map is an
isomorphism and is compatible with every base change.

By Lemma~\ref{lem:G87-projective}, \(P\) is the image of an idempotent
endomorphism of a finite free \(\Lambda_p[\Gamma_q]\)-module.
Apply the functor in (G87.9) and take the image of that idempotent.
Thus \(\mathcal M_\tau\) is a finite projective \(\Lambda_p\)-module,
and (G87.8) follows from the same idempotent formula after arbitrary,
possibly non-flat, base change.
The local ring \(\Lambda_p\) makes this finite projective module free.
A nonzero characteristic-zero fiber then forces its rank to be
positive, so every such fiber is nonzero.
After extending a coefficient field to split the finite group,
a nonzero map from the irreducible \(\tau\) is an occurrence of
\(\tau\). No exactness of the invariants functor on arbitrary
\(\Gamma_q\)-modules is asserted.
\end{proof}

\begin{theorem}[AC-G87: one fixed host for a symmetric weight progression]
\label{thm:G87}
Retain a completed separated instance of AC-G78--AC-G83 with starting
weight \(e\ge1\), including AC-G80's exact old types.
Keep its fields \(F\) and \(H\), its prime \(q\), its character
\(\theta\), all compact levels and projectors, and its residual
away-Hecke maximal ideal \(\mathfrak m\) fixed.
Then for every integer
\[
 w\ge1,\qquad w\equiv e\pmod{268},
 \qquad \mu(w)=(w-1,0,1-w),
 \tag{G87.11}
\]
there is an actual ordinary automorphic representation at weight
\(\mu(w)\) in this same residual block with finite \(q\)-type \(\tau\),
the retained old types, and the exact coefficient-prime flag.
Its normalized rank-three Galois representation is absolutely
irreducible.
After the same fixed soluble base change to \(H\), it remains
absolutely irreducible, is unramified above \(q\), and has the
restrictions of the old types and the corresponding ordinary
coefficient-prime flag.

Thus a single actual \(q\) and a single actual CM cover \(H\) work
for all weights in (G87.11). Their degree and discriminant are the
already obtained fixed quantities of AC-G83 (and AC-G86 if its
explicit least-prime bound is used).
The statement does not prescribe the characteristic-zero Hecke
eigenvalues, the full generic characteristic-\(p\) trace, or a map
from an arbitrary deformation curve to the Hecke algebra.
\end{theorem}

\begin{proof}
\noindent\textbf{The finite torus and center.}
The completed instance supplies a nonzero \(\tau\)-fiber at weight
\(\mu(e)\), with the projectors already imposed.
Let \(\delta_e\) denote its fixed character of the finite Teichmuller
torus.  In the scalar-function convention it is
\[
 \delta_e=
 \left.(\xi_v,1,\xi_v)_v\,
         w_0(\mu(e))^{-1}\right|_{\Delta_p},
 \qquad \xi_v=\chi^j\tau_v .
 \tag{G87.12}
\]
Here the quadratic local character \(\tau_v\), if present in the
original construction, is fixed; it is unrelated to the finite
representation \(\tau=\tau_\theta\).
For parallel algebraic weights, evaluating
\(\mu(w)-\mu(e)\) on a finite residue-field unit factors through
the residue norm to \(\mathbb F_p^\times\).
It is therefore trivial on \(\Delta_p\) when
\(w-e\) is a multiple of \(p-1=268\).
Moreover every \(\mu(w)\) has central exponent
\((w-1)+0+(1-w)=0\).
Thus the fixed finite torus component and the retained central
compatibilities are unchanged; no alteration of a global central
character is being presumed.

\noindent\textbf{Occurrence at every weight.}
Each \(\mu(w)\) is locally dominant algebraic, including \(w=1\).
Proposition~\ref{prop:G87-type} and (G87.4) give a nonzero
\(\tau\)-multiplicity space in the actual
\[
 e_{\rm aux}e_{\rm old}e_{\delta_e}
 H_{\mu(w)}(K^{p,q}K_q^1)_{\mathfrak m}.
 \tag{G87.13}
\]
Thorne's Proposition 6.2(3) identifies this with the corresponding
classical ordinary algebraic modular forms.
The away Hecke algebra commutes with \(\Gamma_q\), so after a finite
coefficient extension there is an eigencharacter in this nonzero
multiplicity space. The classical automorphic decomposition supplies
an actual automorphic constituent containing \(\tau\).
All finite projectors and the localization were present before this
choice, so its old types, finite \(p\)-torus character, and residual
Hecke block are retained.
The coefficient field of this eigensystem is allowed to depend on
\(w\).

\noindent\textbf{The unchanged residual and local arguments.}
Let
\(\bar\gamma_e=\bar\psi\,\bar\epsilon^{\,e}\).
Because \(w\equiv e\pmod{p-1}\), one has
\(\bar\gamma_w=\bar\gamma_e\).
Consequently the normalized residual semisimplification throughout
the progression is the same ordered triple
\[
 (\bar\chi^j\bar\gamma_e)\ \oplus\ 1\
       \oplus\ (\bar\chi^j\bar\gamma_e^{-1}),
 \tag{G87.14}
\]
with the fixed delayed quadratic factor before its killing base
change in the AC-G81 branch.
The source residual genericity witness is unchanged.
Thus its exclusion of noncuspidal discrete blocks and its
split-place genericity conclusion continue to hold.

The exact finite type \(\tau\) gives, by AC-G72-local, the tame
irreducible two-dimensional parameter with inertia
\(\theta\oplus\theta^q\), together with an unramified line, and
monodromy zero.
The normalized residual Frobenius roots at \(q\) remain
\(\{-1,1,-1\}\).
The local two-dimensional block has Frobenius roots \(a,-a\);
hence the unramified line reduces to \(-1\).
In a global conjugate-self-dual decomposition of type \(2+1\),
the unique global line would itself be conjugate self-dual.
The separated triple (G87.14) has only its middle constituent
individually conjugate self-dual, so that global line would reduce
to \(1\), a contradiction at \(q\).
A sum of three global lines cannot contain the irreducible local
two-dimensional block.
The exact AC-G78-type argument therefore proves global absolute
irreducibility for every member of the progression.

At \(p\), the full-torus calculation of AC-G79 recovers the same
smooth finite characters \((\xi_v,1,\xi_v)\) from the fixed
\(\delta_e\) and \(\mu(w)\).
The parallel Hodge exponents are \((-w,0,w)\).
The middle finite inertia character differs from the endpoint
character, while the equal-inertia endpoints have slope gap
\(2wf\), which is never the one-step gap \(f\).
Therefore monodromy is zero, even for \(w=1\).
The same parallel weak-admissibility and descent proof gives
graded characters
\[
 u_{1,v}\chi^j\tau_v\epsilon^{\,w},\qquad
 u_{2,v},\qquad
 u_{3,v}\chi^j\tau_v\epsilon^{-w},
 \tag{G87.15}
\]
where the \(u_{i,v}\) are unramified.
The old-prime rank-one compact types remain imposed in (G87.13);
AC-G80 therefore gives their exact inertia and zero monodromy,
including repeated-character blocks.

\noindent\textbf{The same cover.}
The real cover in AC-G83 was selected to kill the order-\(p\)
\(q\)-inertia \(\theta\oplus\theta^q\).
This type is independent of \(w\), so the same ramification index
\(p\), after the same unramified quadratic local extension, kills it
for every member.
All other assertions about restriction of old inertia and the
coefficient-prime ordinary flag are unchanged.

For completeness, the global irreducibility argument along the
fixed cover uses only rank three and its cyclic prime-degree steps
\(2\) and \(p=269\).
At a cyclic prime-degree step, restriction of an irreducible
characteristic-zero representation is semisimple by the sum of
translates of a simple submodule.
Frobenius reciprocity identifies its endomorphism space with the sum
of the self-twist spaces by characters of that cyclic quotient.
A nontrivial self-twist of a rank-three representation would force
the cube of the twisting character to be trivial, by determinants,
which is impossible for orders \(2\) or \(269\).
Thus restriction stays irreducible at every step.
The same cyclic base-change and definite-unitary descent inputs used
in AC-G83 give the actual representations over \(H\).
All ramified places remain split over the real subfield, and the
coefficient-prime extensions are the same unramified ones.
This proves the theorem with one fixed host.
\end{proof}

\paragraph{Exact scope of weight support.}
The result gives actual eigensystems at each weight in one fixed
residual \emph{maximal-ideal} block.
It does not identify their full characteristic-\(p\) Hecke trace
with the generic trace of an arbitrary characteristic-\(p\)
deformation curve.
Even algebraically, a module may be finite free of positive rank over
a weight ring and still fail to support a prescribed curve in a
larger deformation ring: take
\(\mathcal R=\mathcal O[[W,X]]\), let
\(M=\mathcal O[[W]]\) with \(X\) acting as zero, and use the curve
\(\mathcal R\to k[[t]]\) given by \(W\mapsto t\), \(X\mapsto t\).
This curve is even dominant on the characteristic-\(p\) weight line,
since it identifies \(k[[W]]\) with \(k[[t]]\).
Every weight fiber of \(M\) is nonzero, but its annihilator \((X)\)
is not contained in the curve kernel.
Thus the original per-curve \(J\)-containment and its exact frame
conditions remain separate.

\paragraph{Sources and geometric dictionary.}
The literal source is Thorne,
\emph{A \(p\)-adic approach to the existence of level-raising
congruences}, \S6: Lemma 6.1, the definitions preceding equation
(6.1), that equation and the following depth-stabilization paragraph,
the complete proof of Proposition 6.2, and the finite group-algebra
norm calculation in Proposition 6.7.
These exact applied passages were read in the primary published PDF.
The G78--G83 interfaces, including the full finite-torus normalization,
the rank-one compact types, the finite \(q\)-type, and the fixed cover,
are retained with their proved scopes.
The projective inverse-limit construction and arbitrary-base-change
multiplicity argument above are proved in full.
In the already audited SGA~1 representation/finite-cover dictionary,
one fixed finite cover carries all these restricted local systems.
No existence of a motive, no new archive scan, and no per-curve
automorphy identification follow from this dictionary.
